\section{Conclusions}\label{sec:conclusions}

\paragraph{Association.} In the \eu\ from 2009 to 2026, the negative association between
quarterly GDP growth and the change in the unemployment rate is a robust descriptive
regularity outside the pandemic quarters. A hierarchical model that accounts for
heavy-tailed errors, residual persistence and a common European factor puts the average
country sum at \val{by_mu}, with a between-country standard deviation of \val{by_tau}.

\paragraph{Heterogeneity.} Some country differences are credible in every accepted
specification and persist under other measures of labour input; Spain against the
Netherlands and against Luxembourg are the clearest cases. Most rankings are not
established.

\paragraph{Prediction.} The association adds little to forecasts of unemployment, and
partial pooling does not improve them by more than sampling variation.

\paragraph{Causation.} No identification strategy available with public data turns the
association into a causal effect.

\paragraph{The earlier claims.} Spain's large association is supported. The other two are
qualified. The Dutch delayed response meets its registered rule in the primary fit but not
in every accepted specification. Luxembourg's association is weak and sign-unstable in the
classical estimates, but the registered Bayesian rule for a weak association is met in
none of the accepted fits. The proposed explanation for Luxembourg, the mismatch
between production by cross-border commuters and resident unemployment, is not supported:
domestic and resident employment barely move with GDP, and output per hour absorbs the
fluctuations.

\paragraph{The general lesson.} An Okun coefficient estimated from aggregate data is a
summary of how output and unemployment moved together over a particular period. It is not
a structural parameter of a country's labour market.
